% =====================================================================
% Глава 4. Экспериментальное исследование.
% Версия 2 --- после получения данных E1--E4 + confirmation runs.
% Объем 8--10 страниц.
% =====================================================================

\section{Экспериментальное исследование}
\label{sec:experiments}

Настоящая глава содержит результаты эмпирической проверки
исследовательских вопросов (от англ. Research Questions, RQ)
RQ1--RQ4, поставленных в разделе~\ref{sec:meth:funnel}. Глава
организована как последовательное изложение четырех основных
экспериментов воронки E1--E4, двух подтверждающих прогонов
на~кросс-семейном рабочем агенте и сводного статистического анализа
на~основе непараметрического бутстрапа.

\subsection{Условия проведения и общие характеристики}
\label{sec:exp:setup}

Все эксперименты выполнены в единой программной инфраструктуре,
описанной в главе~\ref{sec:architecture}, в период
с~16 по~17~мая 2026~года. В роли основной языковой модели
(\emph{worker}) для E1--E4 использовалась \code{cerebras:gpt-oss-120b}
с параметром \code{reasoning\_effort=low} для рабочих агентов и
\code{high} для агента-критика. В роли судьи оценки качества
выступала \code{openai:gpt-4.1-mini} с самосогласованностью из трех
независимых вызовов и нулевой температурой выборки.
Шлюз~\code{HumanGateway} (компонент человека в~контуре управления,
от англ. Human-in-the-Loop, HITL) работал в режиме симулятора через
ту же~\code{openai:gpt-4.1-mini} с системной подсказкой, зависящей
от текущей роли. Подтверждающие прогоны (\code{confirmation\_e3},
\code{confirmation\_e4}) выполнены с заменой рабочего агента
на~\code{cerebras:qwen-3-235b-a22b-instruct-2507} (семейство Alibaba
Qwen) при сохранении остального стека неизменным.

Сводные объемные характеристики экспериментальной серии приведены
в~таблице~\ref{tab:exp:totals}. Совокупный объем составил
5175~прогонов и~суммарную стоимость около~\$87~--- существенно ниже
планового бюджета в~\$200~благодаря низкой удельной стоимости
вызовов на~инфраструктуре~Cerebras (от англ. Wafer-Scale Engine,
WSE). Во~всех финальных грид-прогонах доля отказов нулевая.

\begin{table}[!htbp]
\caption{Объемные характеристики экспериментальной серии}
\label{tab:exp:totals}
\small
\begin{tabularx}{\textwidth}{l c c c X}
\toprule
\textbf{Этап} & \textbf{Прогоны} & \textbf{Стоимость, \$} &
\textbf{Время} & \textbf{Назначение} \\
\midrule
E1   & 900  & 12{,}62 & 3{,}7~ч &
   Базовая линия статических топологий, RQ1. \\
E2   & 2295 & 42{,}37 & 13{,}9~ч &
   Влияние роли человека по матрице (топология~$\times$~роль), RQ3. \\
E3   & 540  & 10{,}12 & 12{,}0~ч &
   Адаптивная мета-топология в~трех режимах маршрутизации, RQ2. \\
E4   & 540  & 6{,}38  & 2{,}1~ч &
   Адаптивная роль на~уже-адаптивной топологии, RQ4. \\
\code{conf\_e3} & 360 & 6{,}50  & 1{,}0~ч &
   Кросс-семейная валидация E3 на~семействе Qwen. \\
\code{conf\_e4} & 540 & 9{,}40  & 1{,}3~ч &
   Кросс-семейная валидация E4 на~семействе Qwen. \\
\midrule
\textbf{Итого}  & \textbf{5175}  & \textbf{87{,}39} &
   \textbf{$\approx$34~ч машинного времени} & --- \\
\bottomrule
\end{tabularx}
\end{table}

\subsection{Эксперимент E1 — базовая линия статических топологий}
\label{sec:exp:e1}

Эксперимент~E1 устанавливает верхнюю границу качества для пяти
статических топологий на~четырех корпусах задач~--- программирование
(HumanEval), математические рассуждения (GSM8K), творческая
генерация (CommonGen) и анализ данных (InfiAgent-DABench), далее
DABench,~--- отвечая на~исследовательский вопрос~RQ1. Сетка прогона
составила $5~\text{топологий} \times 4~\text{корпуса} \times
15~\text{задач} \times 3~\text{зерна}~=~900$ прогонов; роль человека
отключена (\code{human.role=none}). Результаты усреднения по~ячейкам
(топология~$\times$~корпус) приведены в~таблице~\ref{tab:exp:e1q}.

\begin{table}[!htbp]
\caption{Среднее качество $q$ статических топологий по корпусам
(жирным выделен победитель в столбце), E1}
\label{tab:exp:e1q}
\small
\begin{tabularx}{\textwidth}{l c c c c c}
\toprule
\textbf{Топология} & \textbf{HumanEval} & \textbf{GSM8K} &
\textbf{CommonGen} & \textbf{DABench} & \textbf{Среднее} \\
\midrule
Chain        & \textbf{0{,}933} & 0{,}911 & 0{,}526 & 0{,}333 & 0{,}676 \\
Star         & 0{,}911 & \textbf{1{,}000} & 0{,}427 & 0{,}282 & 0{,}655 \\
Debate       & 0{,}644 & 0{,}933 & \textbf{0{,}584} & \textbf{0{,}348} & 0{,}627 \\
Hierarchical & 0{,}889 & 0{,}978 & 0{,}486 & 0{,}267 & 0{,}655 \\
Mesh         & 0{,}667 & 0{,}867 & 0{,}397 & 0{,}319 & 0{,}562 \\
\midrule
\textit{Среднее по победителям} & \textit{0{,}933} & \textit{1{,}000}
   & \textit{0{,}584} & \textit{0{,}348} & \textbf{\textit{0{,}716}} \\
\bottomrule
\end{tabularx}
\end{table}

\textbf{Содержательный итог.} Победители на четырех корпусах
различаются~--- chain на HumanEval, star на~GSM8K и debate
на~CommonGen и~DABench. Единой топологии, доминирующей на всех
типах задач, нет. Лучшая глобальная статика~--- chain
со~средним~$q=0{,}676$~--- проигрывает идеальному выбору
\emph{под задачу} (среднее по победителям, $q=0{,}716$) около
четырех процентных пунктов. Этот разрыв задает~\emph{class-level
upper bound} для последующего эксперимента~E3.

Обнаружено два дополнительных явления. Во-первых,
топология~Mesh~деградирует на программировании и творческой
генерации до~$q\le0{,}40$ из-за вырождения голосования
по~совпадению строк. Во-вторых, на корпусе~DABench все статические
топологии демонстрируют~$q\le0{,}35$ при доле полностью неудачных
прогонов~62--71\,\%~--- проявление, которое в~разделе~\ref{sec:exp:conf}
будет переинтерпретировано как специфика рабочего агента,
а не корпуса.

\textbf{Ответ на~RQ1.} Универсального статического победителя
не существует. Оптимальный выбор топологии~--- функция класса
задач, что мотивирует ввод адаптивной мета-топологии в~E3.

\subsection{Эксперимент E2 — роль человека (RQ3)}
\label{sec:exp:e2}

Эксперимент~E2 измеряет влияние роли человека в графе агентов
поверх статических топологий. Сетка построена как декартово
произведение $5~\text{топологий} \times 5~\text{ролей} \times
4~\text{корпуса} \times 15~\text{задач} \times 3~\text{зерна}$
с~двумя фильтрами~--- по~per-task top-3 топологий из~E1 и~по~трем
структурно бессмысленным парам (Debate~$\times$~Coordinator,
Chain~$\times$~Peer, Hierarchical~$\times$~Peer). Итоговый объем
составил~2295~прогонов. Полная матрица победителей роли по~корпусам
приведена в~таблице~\ref{tab:exp:e2roles}.

\begin{table}[!htbp]
\caption{Победитель среди пяти ролей человека по каждому корпусу
с~маргинализацией по~топологиям, E2}
\label{tab:exp:e2roles}
\small
\begin{tabularx}{\textwidth}{l X c c c}
\toprule
\textbf{Корпус} & \textbf{Победившая роль} &
\textbf{$q$ победителя} & \textbf{$n$ ячеек} &
\textbf{$\Delta q$ vs~E1 best static} \\
\midrule
HumanEval & Coordinator & 0{,}948 & 135 & $+0{,}015$ \\
GSM8K     & Monitor     & 0{,}963 & 135 & $-0{,}037$ \\
CommonGen & Peer        & 0{,}643 & 45  & $+0{,}059$ \\
DABench   & Peer        & 0{,}563 & 90  & $+0{,}215$ \\
\bottomrule
\end{tabularx}
\end{table}

Совокупный лифт от~подключения человека~(\emph{HITL lift}),
рассчитанный на~пересечении ячеек~E1~и~E2, составил~$+0{,}033$
по~среднему качеству ($+4{,}8\,\%$ относительного прироста).
Распределение лифта по~ячейкам неоднородно. Положительный лифт
концентрируется на~трех конфигурациях~---
Hierarchical~$\times$~CommonGen ($+0{,}107$),
Chain~$\times$~DABench ($+0{,}087$) и
Mesh~$\times$~DABench ($+0{,}052$); отрицательный лифт наблюдается
только на~Star~$\times$~GSM8K ($-0{,}062$), где топология уже
близка к~потолку качества.

\textbf{Ответ на~RQ3.} Подключение человека дает скромное
улучшение среднего качества при существенно различном эффекте
по~ячейкам. Никакая единая роль не выигрывает на~всех корпусах~---
ключевой аргумент против фиксированного выбора роли и в~пользу
адаптивной маршрутизации, проверяемой в~E4.

\subsection{Эксперимент E3 — адаптивная мета-топология (RQ2)}
\label{sec:exp:e3}

Эксперимент~E3 непосредственно отвечает на~центральный
исследовательский вопрос~RQ2. На~единой адаптивной мета-топологии
сравниваются три режима маршрутизатора~---
\code{rule}~(детерминированное правиловое решение),
\code{llm}~(вызов языковой модели с~парсингом по~Pydantic) и
\code{oracle}~(чтение таблицы~\code{best per task} из~файла,
построенного по~результатам~E1). Каждый режим прогнан в~сетке
$4~\text{корпуса} \times 15~\text{задач} \times 3~\text{зерна}=180$
прогонов; фиксирована роль человека~Reviewer. Сводные показатели
по~режимам приведены в~таблице~\ref{tab:exp:e3overall}.

\begin{table}[!htbp]
\caption{Сводные показатели трех режимов маршрутизатора
топологий, E3}
\label{tab:exp:e3overall}
\small
\begin{tabularx}{\textwidth}{l c c c c X}
\toprule
\textbf{Режим} & \textbf{Среднее $q$} & \textbf{Стоимость, \$} &
\textbf{Время, c} & \textbf{Итераций} & \textbf{Комментарий} \\
\midrule
oracle & 0{,}697 & 0{,}0227 & 1077 & 6{,}27 &
   Верхняя граница, знание корпуса. \\
rule   & 0{,}645 & 0{,}0236 & 928  & 6{,}04 &
   Правиловый, без вызовов языковой модели. \\
llm    & 0{,}631 & \textbf{0{,}0099} & \textbf{881}  & 6{,}10 &
   Вызов языковой модели на каждый переход. \\
\midrule
\textit{Best static (E1)} & \textit{0{,}716} & \textit{0{,}014} &
   \textit{448} & \textit{4{,}35} &
   \textit{Идеальный выбор статики под корпус.} \\
\bottomrule
\end{tabularx}
\end{table}

Численная картина допускает двойственную интерпретацию.
По~абсолютному качеству все три режима маршрутизатора уступают
лучшей статической конфигурации под корпус~--- оракул отстает
на~$1{,}9$ процентных пункта, правиловой~--- на~$7{,}2$,
маршрутизатор на основе языковой модели~--- на~$8{,}5$. Причина~---
организационная накладная стоимость адаптивной обвязки
(фазы~$\times$~подграфы~$\times$~маршрутизатор), которая не
окупается на~хорошо-решаемых корпусах (GSM8K, HumanEval).
Единственный корпус с~положительным абсолютным эффектом адаптации~---
CommonGen, где оракул дает~$+0{,}07$ относительно лучшей статики.

Сравнение трех режимов маршрутизатора между собой~---
\emph{внутри} адаптивного класса~--- выявляет дополнительный
неосновной эффект. Маршрутизатор на~основе языковой модели
проигрывает правиловому всего~$1{,}3$ процентных пункта по~качеству,
но~стоит в~$2{,}38$~раза дешевле. Соотношение стоимостей удерживается
на~всех четырех корпусах в~диапазоне $2{,}12{-}3{,}00$, что
указывает на~\emph{структурный} эффект~--- маршрутизатор на основе
языковой модели систематически выбирает более дешевые подтопологии
(например, Hierarchical с~\code{max\_rounds=2}) против более
прожорливых решений правиловой таблицы. Этот сопутствующий эффект
не отвечает на~RQ2, но представляет самостоятельный интерес и
обсуждается отдельно в~разделе~\ref{sec:exp:conf}.

\textbf{Ответ на~RQ2.} Адаптивная мета-топология не дает значимого
прироста качества и не дает экономии стоимости по~сравнению с~лучшей
статической конфигурацией под корпус. В~качестве референса для
сравнения используется не одна фиксированная статика, а~оракул-выбор
наилучшей статической топологии для каждого корпуса (среднее
качество $0{,}716$). Оракул адаптивной обвязки достигает $0{,}697$
и отстает от~этого референса на~$1{,}9$ процентных пункта; правиловой
и языковой маршрутизаторы~--- на~$7{,}2$ и~$8{,}5$~процентных пункта.
Гипотеза~RQ2 эмпирически не подтверждается. Обнаруженная рядом
сопутствующая находка~--- двукратная экономия стоимости
маршрутизатора на~основе языковой модели внутри адаптивного
класса~--- не отвечает на~RQ2 в~его исходной постановке (сравнение
адаптивной и статической конфигураций), и при кросс-семейной
валидации на~семействе Qwen, где правиловой маршрутизатор Парето-
доминирует над маршрутизатором на~основе языковой модели, эта
находка также не воспроизводится (см.~\ref{sec:exp:conf}).

\subsection{Эксперимент E4 — адаптивная роль (RQ4)}
\label{sec:exp:e4}

Эксперимент~E4 размещает второй слой адаптивности~--- маршрутизатор
ролей~--- поверх уже-адаптивной топологии. На~основе результатов~E3
зафиксирован выбор~\code{topology\_router=llm} как
Парето-оптимальной точки адаптивного класса. Сметаемая переменная~---
\code{role\_router}~$\in$~\{\code{fixed}, \code{rule}, \code{llm}\};
сетка $3 \times 4 \times 15 \times 3=540$ прогонов. Сводка
по~режимам~--- в~таблице~\ref{tab:exp:e4overall}.

\begin{table}[!htbp]
\caption{Сводные показатели трех режимов маршрутизатора ролей
поверх адаптивной мета-топологии, E4}
\label{tab:exp:e4overall}
\small
\begin{tabularx}{\textwidth}{l c c c c}
\toprule
\textbf{Режим} & \textbf{Среднее $q$} & \textbf{$\Delta q$ vs
   \code{fixed}} & \textbf{Стоимость, \$} & \textbf{Время, c} \\
\midrule
fixed & 0{,}653 & ---       & 0{,}01186 & 381 \\
rule  & \textbf{0{,}680} & $+0{,}027$ & 0{,}01193 & 356 \\
llm   & 0{,}643 & $-0{,}010$ & 0{,}01168 & 285 \\
\bottomrule
\end{tabularx}
\end{table}

Маршрутизатор на~основе правил формально увеличивает среднее
качество на~$2{,}7$~процентных пункта над фиксированной ролью
\code{reviewer}. Направление совпадает на~четырех корпусах из~четырех
при~оговорке о~малом размере выборки и одном семействе весов
рабочего агента; драйвер прироста~--- HumanEval ($+0{,}067$).
Стоимость трех режимов укладывается в~коридор шириной~$0{,}85\,\%$~---
дополнительные стоимостные издержки адаптивной маршрутизации роли
незначимы.

Качественный анализ распределения ролей показал, что правиловой
маршрутизатор фактически коллапсирует в~единственную роль~---
координатор активируется в~$98{,}3\,\%$ переходов. То~есть
адаптации роли по~фазам в~текущей реализации не~происходит~---
фактически реализуется почти-статическая политика с~единственной
ролью. E4 переоткрывает результат~E2~\enquote{Coordinator выигрывает
на HumanEval} в~новой обертке. Подробная интерпретация
этого механизма дана в~разделе~\ref{sec:analysis:rq}.

\textbf{Ответ на~RQ4.} Гипотеза о~превосходстве совместной адаптации
топологии и роли над статическими комбинациями эмпирически
не подтверждается. На~исходной семье моделей правиловая
маршрутизация роли дает направленный прирост качества $+0{,}027$
при~нулевой дополнительной стоимости, однако
бутстрап-доверительный интервал $[-0{,}055,\, +0{,}110]$ покрывает
нуль, и~строгий вывод о~значимости прироста не делается. Сопутствующая
находка~--- направление эффекта совпадает на~всех четырех корпусах
из~четырех~--- при~кросс-семейной валидации~(\ref{sec:exp:conf})
также не подтверждается~--- на~семействе Qwen знак эффекта
становится отрицательным.

\subsection{Подтверждающие прогоны — кросс-семейная валидация}
\label{sec:exp:conf}

Два подтверждающих прогона~--- \code{confirmation\_e3}
и~\code{confirmation\_e4}~--- повторяют экспериментальную сетку
E3~и~E4 с~единственной измененной переменной~--- семейством
рабочей языковой модели. Вместо~\code{gpt-oss-120b} использована
\code{qwen-3-235b-a22b-instruct-2507} (семейство Alibaba Qwen).
Все остальные компоненты (судья, шлюз~HITL, инфраструктура,
адаптивная обвязка, защитные правила, конфигурация порогов) сохранены
неизменными. Цель~--- проверить переносимость основных выводов
E3~и~E4 между семействами весов.

Результаты~\code{confirmation\_e3} (\code{topology\_router}~$\in$~\{rule, llm\})
приведены в~таблице~\ref{tab:exp:confe3}. Соотношение стоимостей
правилового маршрутизатора к~маршрутизатору на основе языковой модели
схлопывается с~$2{,}38$ на~исходном семействе до~$1{,}11$~на~семействе
Qwen, а~разрыв качества увеличивается с~$+1{,}3$ процентных пункта
до~$+30{,}6$ в~пользу правилового режима.

\begin{table}[!htbp]
\caption{Сравнение двух режимов маршрутизатора топологий между
семействами весов, \code{confirmation\_e3}}
\label{tab:exp:confe3}
\small
\begin{tabularx}{\textwidth}{l c c c c X}
\toprule
\textbf{Семейство} & \textbf{$q_{\text{rule}}$} &
\textbf{$q_{\text{llm}}$} & \textbf{$\Delta q$} &
\textbf{$c_{\text{rule}}/c_{\text{llm}}$} &
\textbf{Парето-победитель} \\
\midrule
gpt-oss-120b (E3)      & 0{,}645 & 0{,}631 & $+0{,}013$ & $2{,}38$ &
   \code{llm} (по стоимости) \\
qwen-3-235b (conf\_e3) & 0{,}866 & 0{,}560 & \textbf{$+0{,}307$} &
   $1{,}11$ & \code{rule} (по обеим осям) \\
\bottomrule
\end{tabularx}
\end{table}

Аналогичная картина для~\code{confirmation\_e4} (см.~таблицу
\ref{tab:exp:confe4}). На исходном семействе моделей правиловой
маршрутизатор роли превосходил фиксированную конфигурацию
на~$+0{,}027$; на~семействе Qwen он, напротив, уступает
ей~на~$-0{,}081$~--- знак эффекта изменился.

\begin{table}[!htbp]
\caption{Сравнение трех режимов маршрутизатора ролей между
семействами весов, \code{confirmation\_e4}}
\label{tab:exp:confe4}
\small
\begin{tabularx}{\textwidth}{l c c c c}
\toprule
\textbf{Режим} & \textbf{$q$ на gpt-oss} & \textbf{$q$ на qwen} &
\textbf{$\Delta$ vs~\code{fixed} gpt-oss} & \textbf{$\Delta$ vs~\code{fixed} qwen} \\
\midrule
fixed & 0{,}653 & 0{,}582 & ---       & --- \\
rule  & 0{,}680 & 0{,}501 & $+0{,}027$ & \textbf{$-0{,}081$} \\
llm   & 0{,}643 & 0{,}528 & $-0{,}010$ & $-0{,}054$ \\
\bottomrule
\end{tabularx}
\end{table}

Качественный анализ выявил дополнительный эффект~--- сдвиг
распределения ролей при~неизменной таблице~\code{DEFAULT\_ROLE\_TABLE}.
На~исходной семье правиловой маршрутизатор выбирал~\emph{координатор}
в~$98\,\%$ переходов; на~семье~Qwen тот же маршрутизатор с~той же
таблицей выбирает~\emph{равноправного участника} в~$94\,\%$
переходов. Иначе говоря, политика, выглядящая независимой от модели,
фактически зависит от нее через входы~--- фазовые классификаторы
получают разные сигналы от моделей разных семейств и относят те же
ситуации к~разным фазам.

\textbf{Дополнительная находка~--- DABench.} На~исходном семействе
корпус~DABench демонстрировал среднее качество~$0{,}15{-}0{,}24$ во~всех
экспериментах E1--E4; на~семье~Qwen те же конфигурации дают
качество $0{,}84{-}0{,}93$. Это согласуется с~гипотезой о~том, что
DABench измеряет компетенции, специфичные для~семейства весов
рабочего агента, а не является~\enquote{сломанным корпусом};
численный контраст между семействами~--- около 70 процентных пунктов
в~абсолютном выражении~--- исключает корпус как источник сквозного
смещения.

\textbf{Ограничение интерпретации.} На~исходном семействе работники
запускались с~параметром~\code{reasoning\_effort=low}; программный
интерфейс~Qwen этот параметр не~принимает. Часть наблюдаемых
кросс-семейных различий смешана с~различием в~глубине рассуждений
работника и не отделяется от~эффекта семейства весов без
дополнительного контроля; это явное ограничение, требующее
повторной серии при появлении соответствующей опции
в~интерфейсе~Qwen.

\subsection{Бутстрап-доверительные интервалы}
\label{sec:exp:bootstrap}

Для оценки устойчивости заявленных эффектов проведен непараметрический
бутстрап с~десятью тысячами итераций (метод процентилей,
\code{seed=42}). Сводные интервалы для ключевых разностей и
отношений~--- в~таблице~\ref{tab:exp:bootstrap}.

\begin{table}[!htbp]
\caption{Бутстрап-доверительные интервалы $95\,\%$ для ключевых
разностей и отношений}
\label{tab:exp:bootstrap}
\small
\begin{tabularx}{\textwidth}{X c c c}
\toprule
\textbf{Сравнение} & \textbf{Точечная оценка} &
\textbf{Интервал~$95\,\%$} & \textbf{Содержит границу?} \\
\midrule
E3 rule${-}$llm, $q$ (gpt-oss)            & $+0{,}013$ &
   $[-0{,}072,\, +0{,}097]$ & да (нуль) \\
E3 oracle${-}$rule, $q$ (gpt-oss)         & $+0{,}053$ &
   $[-0{,}033,\, +0{,}137]$ & да (нуль) \\
E3 $c_{\text{rule}}/c_{\text{llm}}$       & $2{,}38$   &
   $[1{,}94,\, 2{,}94]$     & \textbf{нет (единицу)} \\
E4 rule${-}$fixed, $q$ (gpt-oss)          & $+0{,}027$ &
   $[-0{,}055,\, +0{,}110]$ & да (нуль) \\
E4 rule${-}$fixed (HumanEval, gpt-oss)    & $+0{,}067$ &
   $[-0{,}022,\, +0{,}156]$ & да (нуль) \\
\code{conf\_e3} rule${-}$llm, $q$ (qwen)  & $+0{,}307$ &
   \textbf{$[+0{,}233,\, +0{,}378]$} & \textbf{нет (нуль)} \\
\code{conf\_e3} $c_{\text{rule}}/c_{\text{llm}}$ (qwen) & $1{,}11$ &
   $[0{,}98,\, 1{,}25]$     & да (единицу) \\
\code{conf\_e4} rule${-}$fixed, $q$ (qwen) & $-0{,}081$ &
   $[-0{,}173,\, +0{,}010]$ & пограничное \\
\bottomrule
\end{tabularx}
\end{table}

Интервалы позволяют отделить устойчивые эффекты от направленных
точечных оценок. Устойчивых эффектов в~работе три. Во-первых,
двукратная экономия стоимости маршрутизатора на основе языковой
модели на~исходной семье~--- интервал отношения стоимостей
$[1{,}94,\, 2{,}94]$ никогда не пересекает~единицу. Во-вторых,
изменение знака эффекта качества правилового маршрутизатора при
смене семьи весов~--- интервал разности качеств на~семье~Qwen
$[+0{,}233,\, +0{,}378]$ устойчиво исключает нуль. В-третьих,
схлопывание соотношения стоимостей при смене семьи~--- интервал
$[0{,}98,\, 1{,}25]$ покрывает единицу, что указывает на
структурный характер схлопывания. Остальные точечные оценки
направленные, но не устойчивы~--- бутстрап-интервалы включают
нуль, и эффекты в строгом смысле статистически неотличимы от
нулевого.

\subsection{Абляционное исследование}
\label{sec:exp:ablation}

Поскольку структура воронки от~E1 до~E4 уже устроена как
последовательное~\emph{добавление} компонентов~--- сначала статические
топологии, затем человек, затем адаптивная топология, затем
адаптивная роль~--- сама воронка играет роль абляционного исследования.
Сводка эффектов компонентов приведена в~таблице~\ref{tab:exp:ablation}.

\begin{table}[!htbp]
\caption{Абляционная сводка~--- эффект добавления компонента
относительно базовой линии}
\label{tab:exp:ablation}
\small
\begin{tabularx}{\textwidth}{l c c X}
\toprule
\textbf{Компонент} & \textbf{$\Delta q$} & \textbf{$\Delta c$} &
\textbf{Источник эффекта} \\
\midrule
Best static per task (E1)            & ---       & ---       &
   Базовая линия. \\
+ HITL фиксированной роли (E2)       & $+0{,}033$ & $+0{,}004$ &
   Подключение симулятора человека. \\
+ Адаптивная мета-топология, oracle (E3)
                                     & $-0{,}019$ & $+0{,}008$ &
   Накладные расходы маршрутизации без явного выигрыша. \\
+ Адаптивная роль, rule (E4 vs E3 llm)
                                     & $+0{,}049$ & $\approx 0$ &
   Поверх E3 llm; скрытое смещение к роли~\emph{coordinator}. \\
+ Смена семейства модели (\code{conf\_e4} vs E4)
                                     & $-0{,}108$ & $+0{,}005$ &
   Падение по знаку и магнитуде~--- переносимость не подтверждена. \\
\bottomrule
\end{tabularx}
\end{table}

Содержательно абляция показывает, что только два компонента~---
подключение человека (E2~vs~E1) и фиксация роли~\emph{coordinator}
через правиловую таблицу (E4~vs~E3)~--- дают положительное
приращение качества; адаптивная топология сама по себе ухудшает
качество относительно идеально-настроенной статики; смена семейства
весов меняет знак эффекта роли. Кросс-семейная неустойчивость
ставит под вопрос интерпретацию положительных приращений как
универсальных, а не специфичных для исходной семьи моделей.

\subsection{Выводы по главе}
\label{sec:exp:summary}

Результаты экспериментальной серии сводятся к~четырем содержательным
итогам. Первый~--- класс задач определяет оптимальную статическую
топологию (RQ1~подтверждается). Второй~--- адаптивная мета-топология
не превосходит лучшую статическую конфигурацию под корпус
ни~по~качеству, ни~по~стоимости (RQ2~не~подтверждается); рядом
обнаружен неосновной структурный эффект~--- двукратная экономия
стоимости маршрутизатора на~основе языковой модели внутри
адаптивного класса~--- который при~кросс-семейной валидации также
не подтверждается. Третий~--- подключение человека дает скромный
направленный прирост качества при~существенно разном эффекте
по~ячейкам (RQ3~подтверждается). Четвертый~--- совместная адаптация
топологии и роли не превосходит фиксированные комбинации
(RQ4~не~подтверждается); направление эффекта совпадает
на~четырех корпусах из~четырех в~исходной семье моделей, но при
кросс-семейной валидации знак меняется и эффект исчезает.

Единственный устойчивый кросс-семейный вывод~--- зависимость
оптимальной политики маршрутизации от~семейства весов рабочего
агента. Его содержательная интерпретация и обсуждение угроз
валидности вынесены в~главу~\ref{sec:analysis}.
